% ch0-preliminaires.tex — Chapitre 0 : Préliminaires
% Ce chapitre rappelle les notions fondamentales de logique, de théorie
% des ensembles et de structures algébriques nécessaires à l'étude de
% l'algèbre linéaire.

\chapter{Préliminaires}\label{ch:preliminaires}

\section*{Introduction}
\addcontentsline{toc}{section}{Introduction}

L'algèbre linéaire repose sur un socle de notions que l'on suppose
généralement acquises~: ensembles, applications, relations, et surtout
les structures algébriques élémentaires que sont les groupes, les anneaux
et les corps. Ce chapitre zéro a pour vocation de \emph{rappeler} et de
\emph{préciser} ces outils, afin que le lecteur dispose d'un langage
commun et rigoureux pour la suite du cours.

Nous y revisiterons également les principales techniques de démonstration
--- preuve directe, contraposée, raisonnement par l'absurde, récurrence ---
qui seront utilisées de manière intensive dans les chapitres suivants.

Le lecteur déjà à l'aise avec ces prérequis pourra parcourir ce chapitre
rapidement, en s'attardant sur les exercices pour vérifier sa maîtrise.

% ============================================================
\section{Ensembles}\label{sec:ensembles}
% ============================================================

\begin{definition}{Ensemble}{def:ensemble}
  Un \emph{ensemble}\index{ensemble} est une collection d'objets, appelés
  \emph{éléments}, réunis selon un critère bien défini. On note $x\in E$
  le fait que $x$ est un élément de $E$, et $x\notin E$ le contraire.
\end{definition}

On rappelle les ensembles de nombres classiques~:
\[
  \N \subset \Z \subset \Q \subset \R \subset \C.
\]

\begin{definition}{Opérations ensemblistes}{def:operations-ensemblistes}
  Soient $A$ et $B$ deux sous-ensembles d'un ensemble~$E$.
  \begin{enumerate}[label=(\roman*)]
    \item \textbf{Réunion~:} $A\cup B \deff \setcond{x\in E}{x\in A\text{ ou }x\in B}$.
    \item \textbf{Intersection~:} $A\cap B \deff \setcond{x\in E}{x\in A\text{ et }x\in B}$.
    \item \textbf{Différence~:} $A\setminus B \deff \setcond{x\in E}{x\in A\text{ et }x\notin B}$.
    \item \textbf{Complémentaire~:} $A^c \deff E\setminus A$.
    \item \textbf{Produit cartésien~:} $A\times B \deff \setcond{(a,b)}{a\in A,\; b\in B}$.
    \item \textbf{Ensemble des parties~:} $\mathcal{P}(E) \deff \setcond{A}{A\subseteq E}$.
  \end{enumerate}
\end{definition}

\begin{example}{Parties de $\set{1,2,3}$}{ex:parties}
  Soit $E=\set{1,2,3}$. Alors
  \[
    \mathcal{P}(E)=\set{\emptyset,\;\set{1},\;\set{2},\;\set{3},\;
    \set{1,2},\;\set{1,3},\;\set{2,3},\;\set{1,2,3}}.
  \]
  On a $\card{\mathcal{P}(E)}=2^3=8$. Plus généralement, si $E$ est fini
  de cardinal~$n$, alors $\card{\mathcal{P}(E)}=2^n$.
\end{example}

\begin{proposition}{Lois de De Morgan}{prop:de-morgan}
  Soient $A,B\subseteq E$. Alors~:
  \[
    (A\cup B)^c = A^c \cap B^c
    \qquad\text{et}\qquad
    (A\cap B)^c = A^c \cup B^c.
  \]
\end{proposition}

\begin{proof}
  Montrons la première égalité. Soit $x\in E$.
  \begin{align*}
    x\in (A\cup B)^c
      &\iff x\notin A\cup B \\
      &\iff x\notin A \text{ et } x\notin B \\
      &\iff x\in A^c \text{ et } x\in B^c \\
      &\iff x\in A^c\cap B^c.
  \end{align*}
  La seconde se démontre de manière analogue (ou par application de la
  première au complémentaire).
\end{proof}

% ============================================================
\section{Applications}\label{sec:applications}
% ============================================================

\begin{definition}{Application}{def:application}
  Soient $E$ et $F$ deux ensembles. Une \emph{application}\index{application}
  (ou \emph{fonction}) de $E$ dans $F$ est une règle $f$ qui, à chaque
  élément $x\in E$, associe un unique élément $f(x)\in F$. On note
  \[
    f\colon E\to F,\qquad x\mapsto f(x).
  \]
  L'ensemble $E$ s'appelle l'\emph{ensemble de départ} (ou \emph{domaine}),
  et $F$ l'\emph{ensemble d'arrivée} (ou \emph{codomaine}).
\end{definition}

\begin{definition}{Image et image réciproque}{def:image}
  Soit $f\colon E\to F$ une application.
  \begin{enumerate}[label=(\roman*)]
    \item L'\emph{image directe} d'une partie $A\subseteq E$ est
      $f(A)\deff\setcond{f(x)}{x\in A}$.
    \item L'\emph{image réciproque} d'une partie $B\subseteq F$ est
      $f^{-1}(B)\deff\setcond{x\in E}{f(x)\in B}$.
    \item L'\emph{image} de $f$ est $\im f\deff f(E)$.
  \end{enumerate}
\end{definition}

\begin{definition}{Injectivité, surjectivité, bijectivité}{def:inj-surj-bij}
  Soit $f\colon E\to F$ une application.
  \begin{enumerate}[label=(\roman*)]
    \item $f$ est \emph{injective}\index{injection} si
      $\forall\, x_1,x_2\in E,\; f(x_1)=f(x_2)\implies x_1=x_2$.
    \item $f$ est \emph{surjective}\index{surjection} si
      $\forall\, y\in F,\;\exists\, x\in E,\; f(x)=y$
      (autrement dit, $\im f = F$).
    \item $f$ est \emph{bijective}\index{bijection} si elle est à la fois
      injective et surjective.
  \end{enumerate}
\end{definition}

\begin{example}{Exemples fondamentaux}{ex:inj-surj}
  \begin{enumerate}[label=(\alph*)]
    \item L'application $f\colon\R\to\R$, $x\mapsto x^2$ n'est ni
      injective (car $f(-1)=f(1)=1$) ni surjective (car $-1\notin\im f$).
    \item L'application $g\colon\R\to\R_+$, $x\mapsto x^2$ est
      surjective mais pas injective.
    \item L'application $h\colon\R_+\to\R_+$, $x\mapsto x^2$ est bijective,
      de réciproque $h^{-1}\colon y\mapsto\sqrt{y}$.
  \end{enumerate}
\end{example}

\begin{definition}{Composition}{def:composition}
  Soient $f\colon E\to F$ et $g\colon F\to G$ deux applications.
  La \emph{composée}\index{composition} de $f$ et $g$ est l'application
  \[
    g\circ f\colon E\to G,\qquad x\mapsto g(f(x)).
  \]
\end{definition}

\begin{proposition}{Composition et injectivité/surjectivité}{prop:comp-inj-surj}
  Soient $f\colon E\to F$ et $g\colon F\to G$.
  \begin{enumerate}[label=(\roman*)]
    \item Si $g\circ f$ est injective, alors $f$ est injective.
    \item Si $g\circ f$ est surjective, alors $g$ est surjective.
    \item $g\circ f$ est bijective si et seulement si $f$ est injective
      et $g$ est surjective \emph{et} $\im f = F$ (c'est-à-dire $f$ bijective
      et $g$ bijective).
  \end{enumerate}
\end{proposition}

\begin{proof}
  \emph{(i)} Supposons $g\circ f$ injective. Soient $x_1,x_2\in E$ tels
  que $f(x_1)=f(x_2)$. Alors $g(f(x_1))=g(f(x_2))$, donc
  $(g\circ f)(x_1)=(g\circ f)(x_2)$, d'où $x_1=x_2$ par injectivité de
  $g\circ f$.

  \emph{(ii)} Supposons $g\circ f$ surjective. Soit $z\in G$. Il existe
  $x\in E$ tel que $g(f(x))=z$. En posant $y=f(x)\in F$, on obtient
  $g(y)=z$, donc $g$ est surjective.

  Le point \emph{(iii)} se déduit des deux premiers (exercice).
\end{proof}

% ============================================================
\section{Relations}\label{sec:relations}
% ============================================================

\begin{definition}{Relation binaire}{def:relation}
  Une \emph{relation binaire}\index{relation binaire} sur un ensemble $E$
  est un sous-ensemble $\mathcal{R}\subseteq E\times E$. On note
  $x\,\mathcal{R}\,y$ pour $(x,y)\in\mathcal{R}$.
\end{definition}

\subsection{Relations d'équivalence}\label{subsec:equivalence}

\begin{definition}{Relation d'équivalence}{def:rel-equivalence}
  Une relation $\sim$ sur $E$ est une \emph{relation d'équivalence}%
  \index{relation d'équivalence} si elle vérifie~:
  \begin{enumerate}[label=(\roman*)]
    \item \textbf{Réflexivité~:} $\forall\, x\in E,\; x\sim x$.
    \item \textbf{Symétrie~:} $\forall\, x,y\in E,\; x\sim y\implies y\sim x$.
    \item \textbf{Transitivité~:} $\forall\, x,y,z\in E,\;
      (x\sim y\text{ et }y\sim z)\implies x\sim z$.
  \end{enumerate}
  La \emph{classe d'équivalence} de $x$ est
  $\bar{x}\deff\setcond{y\in E}{y\sim x}$.
  L'\emph{ensemble quotient} est $E/{\sim}\;\deff\setcond{\bar{x}}{x\in E}$.
\end{definition}

\begin{example}{Congruence modulo $n$}{ex:congruence}
  Sur $\Z$, la relation définie par
  \[
    a\equiv b\pmod{n} \iff n\mid(a-b)
  \]
  est une relation d'équivalence. L'ensemble quotient est
  $\Z/n\Z = \set{\bar{0},\bar{1},\ldots,\overline{n-1}}$, l'ensemble
  des classes de congruence modulo~$n$. Par exemple, pour $n=3$~:
  \[
    \bar{0}=\set{\ldots,-6,-3,0,3,6,\ldots},\quad
    \bar{1}=\set{\ldots,-5,-2,1,4,7,\ldots},\quad
    \bar{2}=\set{\ldots,-4,-1,2,5,8,\ldots}.
  \]
\end{example}

\subsection{Relations d'ordre}\label{subsec:ordre}

\begin{definition}{Relation d'ordre}{def:rel-ordre}
  Une relation $\leq$ sur $E$ est une \emph{relation d'ordre}%
  \index{relation d'ordre} si elle vérifie~:
  \begin{enumerate}[label=(\roman*)]
    \item \textbf{Réflexivité~:} $\forall\, x\in E,\; x\leq x$.
    \item \textbf{Antisymétrie~:} $\forall\, x,y\in E,\;
      (x\leq y\text{ et }y\leq x)\implies x=y$.
    \item \textbf{Transitivité~:} $\forall\, x,y,z\in E,\;
      (x\leq y\text{ et }y\leq z)\implies x\leq z$.
  \end{enumerate}
  L'ordre est \emph{total} si $\forall\, x,y\in E$, on a $x\leq y$ ou
  $y\leq x$~; sinon il est \emph{partiel}.
\end{definition}

\begin{example}{Ordres classiques}{ex:ordres}
  \begin{enumerate}[label=(\alph*)]
    \item L'ordre usuel $\leq$ sur $\R$ est un ordre total.
    \item L'inclusion $\subseteq$ sur $\mathcal{P}(E)$ est un ordre
      partiel (dès que $\card{E}\geq 2$, il existe des parties
      incomparables).
    \item La divisibilité $a\mid b$ sur $\N^*$ est un ordre partiel.
  \end{enumerate}
\end{example}

% ============================================================
\section{Structures algébriques}\label{sec:structures}
% ============================================================

Cette section introduit les trois structures fondamentales --- groupe,
anneau, corps --- en insistant sur les axiomes et les exemples concrets.
En algèbre linéaire, le corps de base (souvent $\R$ ou $\C$) joue un
rôle central~; il est donc essentiel d'en comprendre les propriétés.

\subsection{Groupes}\label{subsec:groupes}

\begin{definition}{Loi de composition interne}{def:lci}
  Une \emph{loi de composition interne}\index{loi de composition interne}
  sur un ensemble $E$ est une application $\star\colon E\times E\to E$.
  On note $x\star y$ l'image du couple $(x,y)$.
\end{definition}

\begin{definition}{Groupe}{def:groupe}
  Un \emph{groupe}\index{groupe} est un couple $(G,\star)$ où $G$ est un
  ensemble non vide et $\star$ est une loi de composition interne sur $G$
  vérifiant~:
  \begin{enumerate}[label=\textbf{(G\arabic*)}]
    \item \textbf{Associativité~:} $\forall\, a,b,c\in G,\;
      (a\star b)\star c = a\star(b\star c)$.
    \item \textbf{Élément neutre~:} $\exists\, e\in G,\;
      \forall\, a\in G,\; a\star e = e\star a = a$.
    \item \textbf{Symétrique~:} $\forall\, a\in G,\;\exists\, a'\in G,\;
      a\star a' = a'\star a = e$.
  \end{enumerate}
  Le groupe est dit \emph{abélien}\index{groupe!abélien} (ou
  \emph{commutatif}) si de plus~:
  \begin{enumerate}[label=\textbf{(G\arabic*)}]
    \setcounter{enumi}{3}
    \item \textbf{Commutativité~:} $\forall\, a,b\in G,\;
      a\star b = b\star a$.
  \end{enumerate}
\end{definition}

\begin{example}{Groupes classiques}{ex:groupes}
  \begin{enumerate}[label=(\alph*)]
    \item $(\Z,+)$ est un groupe abélien, d'élément neutre $0$.
      Le symétrique de $a$ est $-a$.
    \item $(\Q^*,\times)$ est un groupe abélien, d'élément neutre $1$.
      Le symétrique de $a\neq 0$ est $1/a$.
    \item $(\mathfrak{S}_n,\circ)$, le groupe des permutations de
      $\set{1,\ldots,n}$, est un groupe non abélien dès que $n\geq 3$,
      d'ordre $n!$.
    \item $(\GL_n(\R),\times)$, le groupe des matrices inversibles
      $n\times n$ à coefficients réels, est un groupe non abélien dès
      que $n\geq 2$.
  \end{enumerate}
\end{example}

\begin{proposition}{Unicité du neutre et des symétriques}{prop:unicite-neutre}
  Dans un groupe $(G,\star)$~:
  \begin{enumerate}[label=(\roman*)]
    \item l'élément neutre est unique~;
    \item pour tout $a\in G$, le symétrique de $a$ est unique.
  \end{enumerate}
\end{proposition}

\begin{proof}
  \emph{(i)} Soient $e$ et $e'$ deux éléments neutres. Alors
  $e = e\star e' = e'$, puisque $e'$ est neutre à droite et $e$ est
  neutre à gauche.

  \emph{(ii)} Soient $a'$ et $a''$ deux symétriques de $a$. Alors
  $a' = a'\star e = a'\star(a\star a'') = (a'\star a)\star a'' = e\star a'' = a''$.
\end{proof}

\subsection{Anneaux}\label{subsec:anneaux}

\begin{definition}{Anneau}{def:anneau}
  Un \emph{anneau}\index{anneau} est un triplet $(A,+,\times)$ où $A$ est
  un ensemble muni de deux lois de composition internes vérifiant~:
  \begin{enumerate}[label=\textbf{(A\arabic*)}]
    \item $(A,+)$ est un groupe abélien (de neutre $0_A$).
    \item \textbf{Associativité de $\times$~:} $\forall\, a,b,c\in A,\;
      (a\times b)\times c = a\times(b\times c)$.
    \item \textbf{Élément unité~:} $\exists\, 1_A\in A,\;
      \forall\, a\in A,\; 1_A\times a = a\times 1_A = a$.
    \item \textbf{Distributivité~:} $\forall\, a,b,c\in A$,
      \[
        a\times(b+c)=a\times b + a\times c
        \qquad\text{et}\qquad
        (a+b)\times c = a\times c + b\times c.
      \]
  \end{enumerate}
  L'anneau est \emph{commutatif}\index{anneau!commutatif} si de plus
  $a\times b = b\times a$ pour tous $a,b\in A$.
\end{definition}

\begin{remark}{Conventions}{rem:anneau-convention}
  Certains auteurs n'exigent pas l'existence d'un élément unité~; on
  parle alors d'\emph{anneau unitaire} lorsque $1_A$ existe. Dans ce
  cours, tous nos anneaux seront unitaires.
\end{remark}

\begin{example}{Anneaux classiques}{ex:anneaux}
  \begin{enumerate}[label=(\alph*)]
    \item $(\Z,+,\times)$ est un anneau commutatif.
    \item $(\Mat_n(\R),+,\times)$ est un anneau non commutatif pour $n\geq 2$.
    \item $(\Z/n\Z,+,\times)$ est un anneau commutatif.
    \item L'anneau des polynômes $\R[X]$ est un anneau commutatif.
  \end{enumerate}
\end{example}

\subsection{Corps}\label{subsec:corps}

\begin{definition}{Corps}{def:corps}
  Un \emph{corps}\index{corps} est un triplet $(\K,+,\times)$ tel que~:
  \begin{enumerate}[label=\textbf{(C\arabic*)}]
    \item $(\K,+,\times)$ est un anneau commutatif.
    \item $\K$ possède au moins deux éléments~: $0_\K\neq 1_\K$.
    \item \textbf{Inverse multiplicatif~:} $\forall\, a\in\K\setminus\set{0_\K},\;
      \exists\, a^{-1}\in\K,\; a\times a^{-1}=1_\K$.
  \end{enumerate}
  Autrement dit, $(\K\setminus\set{0_\K},\times)$ est un groupe abélien.
\end{definition}

\begin{remark}{Corps et algèbre linéaire}{rem:corps-al}
  Dans ce cours, la lettre $\K$ désignera toujours un corps, le plus
  souvent $\R$ ou $\C$. Tous les espaces vectoriels, toutes les matrices,
  toutes les applications linéaires seront définis \emph{sur un corps}~$\K$.
  La compréhension des axiomes de corps est donc indispensable.
\end{remark}

% ============================================================
\section{Les axiomes de corps en détail}\label{sec:axiomes-corps}
% ============================================================

Récapitulons de manière systématique. Un corps $(\K,+,\times)$ satisfait
les axiomes suivants, pour tous $a,b,c\in\K$~:

\bigskip
\begin{center}
\renewcommand{\arraystretch}{1.4}
\begin{tabular}{@{}lll@{}}
\toprule
\textbf{Axiome} & \textbf{Addition} & \textbf{Multiplication} \\
\midrule
Associativité   & $(a+b)+c=a+(b+c)$  & $(a\times b)\times c = a\times(b\times c)$ \\
Commutativité   & $a+b=b+a$          & $a\times b = b\times a$ \\
Élément neutre  & $a+0=a$            & $a\times 1=a$ \\
Inverse         & $a+(-a)=0$         & $a\times a^{-1}=1$ \; ($a\neq 0$) \\
Distributivité  & \multicolumn{2}{l}{$a\times(b+c)=a\times b + a\times c$} \\
Non-trivialité  & \multicolumn{2}{l}{$0\neq 1$} \\
\bottomrule
\end{tabular}
\end{center}

\begin{proposition}{Propriétés élémentaires d'un corps}{prop:corps-elem}
  Dans un corps $(\K,+,\times)$~:
  \begin{enumerate}[label=(\roman*)]
    \item $\forall\, a\in\K,\; 0\times a = 0$.
    \item $\forall\, a,b\in\K,\; (-a)\times b = -(a\times b) = a\times(-b)$.
    \item $\forall\, a,b\in\K,\; a\times b = 0 \implies a=0\text{ ou }b=0$
      \emph{(intégrité)}.
  \end{enumerate}
\end{proposition}

\begin{proof}
  \emph{(i)} On a $0\times a = (0+0)\times a = 0\times a + 0\times a$.
  En ajoutant $-(0\times a)$ des deux côtés, on obtient $0=0\times a$.

  \emph{(ii)} On a $a\times b + (-a)\times b = (a+(-a))\times b = 0\times b = 0$,
  donc $(-a)\times b = -(a\times b)$.

  \emph{(iii)} Supposons $a\neq 0$. Alors $a^{-1}$ existe, et
  $b = 1\times b = (a^{-1}\times a)\times b = a^{-1}\times(a\times b) =
  a^{-1}\times 0 = 0$.
\end{proof}

\subsection{Le corps $\R$ des nombres réels}\label{subsec:corps-R}

\begin{example}{Le corps $\R$}{ex:corps-R}
  L'ensemble $\R$ des nombres réels, muni de l'addition et de la
  multiplication usuelles, forme un corps. C'est le corps le plus
  familier~; il possède en outre une structure d'ordre compatible~: pour
  tous $a,b,c\in\R$,
  \[
    a\leq b \implies a+c\leq b+c,
    \qquad
    (a\leq b\text{ et }c\geq 0) \implies ac\leq bc.
  \]
  On dit que $\R$ est un \emph{corps ordonné}. De plus, $\R$ vérifie
  l'\emph{axiome de la borne supérieure}~: toute partie non vide et
  majorée de $\R$ admet une borne supérieure. Cette propriété le
  distingue de $\Q$ et sera cruciale en analyse.
\end{example}

\subsection{Le corps $\C$ des nombres complexes}\label{subsec:corps-C}

\begin{definition}{Nombres complexes}{def:complexes}
  L'ensemble des \emph{nombres complexes}\index{nombres complexes} est
  \[
    \C \deff \setcond{a+bi}{a,b\in\R},
  \]
  où $i$ est un symbole vérifiant $i^2=-1$. Les opérations sont~:
  \begin{align*}
    (a+bi)+(c+di) &= (a+c)+(b+d)i, \\
    (a+bi)\times(c+di) &= (ac-bd)+(ad+bc)i.
  \end{align*}
\end{definition}

\begin{example}{Inverse dans $\C$}{ex:inverse-C}
  Soit $z=2+3i\in\C$, $z\neq 0$. Son inverse est
  \[
    z^{-1} = \frac{1}{2+3i} = \frac{2-3i}{(2+3i)(2-3i)}
    = \frac{2-3i}{4+9} = \frac{2}{13}-\frac{3}{13}i.
  \]
  On utilise la formule générale~: si $z=a+bi\neq 0$, alors
  $z^{-1} = \frac{a-bi}{a^2+b^2}$.
\end{example}

\begin{remark}{$\C$ n'est pas ordonné}{rem:C-non-ordonne}
  Contrairement à $\R$, le corps $\C$ ne peut pas être muni d'un ordre
  total compatible avec les opérations. En revanche, $\C$ est
  \emph{algébriquement clos}~: tout polynôme non constant à coefficients
  complexes admet au moins une racine dans $\C$. Ce résultat fondamental,
  le \emph{théorème de d'Alembert--Gauss}, sera au c\oe{}ur de la théorie
  spectrale (
ef{ch:preliminaires} et chapitres ultérieurs).
\end{remark}

\subsection{Les corps finis $\F_p$}\label{subsec:corps-Fp}

\begin{definition}{Corps fini $\F_p$}{def:Fp}
  Soit $p$ un nombre premier. Le \emph{corps fini}\index{corps fini}
  $\F_p$ est l'ensemble $\Z/p\Z = \set{\bar{0},\bar{1},\ldots,\overline{p-1}}$
  muni de l'addition et de la multiplication modulo~$p$.
\end{definition}

\begin{proposition}{$\Z/p\Z$ est un corps si et seulement si $p$ est premier}{prop:Zp-corps}
  L'anneau $(\Z/n\Z,+,\times)$ est un corps si et seulement si $n$ est
  un nombre premier.
\end{proposition}

\begin{proof}
  $(\Rightarrow)$ Si $n$ n'est pas premier, écrivons $n=ab$ avec
  $1<a,b<n$. Alors $\bar{a}\neq\bar{0}$ et $\bar{b}\neq\bar{0}$, mais
  $\bar{a}\times\bar{b}=\overline{ab}=\bar{0}$. L'anneau possède des
  diviseurs de zéro, donc n'est pas un corps.

  $(\Leftarrow)$ Si $p$ est premier, soit $\bar{a}\in\Z/p\Z$ avec
  $\bar{a}\neq\bar{0}$, c'est-à-dire $p\nmid a$. Par le théorème de
  Bézout, il existe $u,v\in\Z$ tels que $au+pv=1$, d'où
  $\bar{a}\times\bar{u}=\bar{1}$. Ainsi $\bar{a}$ est inversible.
\end{proof}

\begin{example}{Tables d'opérations de $\F_3$}{ex:F3}
  Dans $\F_3=\set{0,1,2}$~:

  \medskip
  \begin{minipage}{0.45\textwidth}
  \centering
  \begin{tabular}{c|ccc}
    $+$ & $0$ & $1$ & $2$ \\\hline
    $0$ & $0$ & $1$ & $2$ \\
    $1$ & $1$ & $2$ & $0$ \\
    $2$ & $2$ & $0$ & $1$ \\
  \end{tabular}
  \end{minipage}
  \begin{minipage}{0.45\textwidth}
  \centering
  \begin{tabular}{c|ccc}
    $\times$ & $0$ & $1$ & $2$ \\\hline
    $0$ & $0$ & $0$ & $0$ \\
    $1$ & $0$ & $1$ & $2$ \\
    $2$ & $0$ & $2$ & $1$ \\
  \end{tabular}
  \end{minipage}

  \medskip
  On vérifie que $1^{-1}=1$ et $2^{-1}=2$ (car $2\times 2=4\equiv 1\pmod{3}$).
\end{example}

\begin{example}{Le corps $\F_5$}{ex:F5}
  Dans $\F_5=\set{0,1,2,3,4}$, calculons les inverses multiplicatifs~:
  \[
    1^{-1}=1,\quad 2^{-1}=3\;\;(2\times 3=6\equiv 1),\quad
    3^{-1}=2,\quad 4^{-1}=4\;\;(4\times 4=16\equiv 1).
  \]
  On peut faire de l'algèbre linéaire sur $\F_5$~: par exemple, résoudre
  un système d'équations linéaires à coefficients dans $\F_5$.
\end{example}

\begin{remark}{Caractéristique d'un corps}{rem:caracteristique}
  La \emph{caractéristique}\index{caractéristique} d'un corps $\K$,
  notée $\car(\K)$, est le plus petit entier $p>0$ tel que
  $\underbrace{1_\K+\cdots+1_\K}_{p}=0_\K$, ou $0$ si un tel entier
  n'existe pas. On a $\car(\R)=\car(\C)=\car(\Q)=0$ et $\car(\F_p)=p$.
  Lorsque $\car(\K)>0$, elle est nécessairement un nombre premier.
\end{remark}

% ============================================================
\section{Techniques de démonstration}\label{sec:demonstrations}
% ============================================================

Les démonstrations sont le c\oe{}ur des mathématiques. Dans ce cours,
nous rédigerons des preuves complètes et nous attendons du lecteur qu'il
en fasse autant dans les exercices. Rappelons les techniques les plus
courantes.

\subsection{Preuve directe}\label{subsec:preuve-directe}

\begin{remark}{Méthode}{rem:preuve-directe}
  Pour montrer $P\implies Q$, on suppose $P$ vraie et on en déduit $Q$
  par une chaîne d'implications logiques.
\end{remark}

\begin{example}{Preuve directe}{ex:preuve-directe}
  \textbf{Énoncé.} Si $n$ est pair, alors $n^2$ est pair.

  \textbf{Preuve.} Supposons $n$ pair~: il existe $k\in\Z$ tel que $n=2k$.
  Alors $n^2=(2k)^2=4k^2=2(2k^2)$. Comme $2k^2\in\Z$, on a que $n^2$ est
  pair.
\end{example}

\subsection{Preuve par contraposée}\label{subsec:contraposee}

\begin{remark}{Méthode}{rem:contraposee}
  L'implication $P\implies Q$ est logiquement équivalente à sa
  \emph{contraposée} $\lnot Q\implies\lnot P$. On peut donc montrer
  $P\implies Q$ en supposant $\lnot Q$ et en déduisant $\lnot P$.
\end{remark}

\begin{example}{Preuve par contraposée}{ex:contraposee}
  \textbf{Énoncé.} Si $n^2$ est impair, alors $n$ est impair.

  \textbf{Preuve (par contraposée).} On montre~: si $n$ est pair, alors
  $n^2$ est pair. Or c'est exactement le résultat de
  l'
ef{ex:ex:preuve-directe}. \qed
\end{example}

\subsection{Preuve par l'absurde}\label{subsec:absurde}

\begin{remark}{Méthode}{rem:absurde}
  Pour montrer une proposition $P$, on suppose $\lnot P$ et on en déduit
  une contradiction. On conclut alors que $P$ est vraie.
\end{remark}

\begin{example}{$\sqrt{2}$ est irrationnel}{ex:sqrt2}
  \textbf{Énoncé.} $\sqrt{2}\notin\Q$.

  \textbf{Preuve.} Supposons par l'absurde que $\sqrt{2}\in\Q$. Il
  existe alors $p,q\in\Z$, $q\neq 0$, avec $\pgcd(p,q)=1$, tels que
  $\sqrt{2}=p/q$. En élevant au carré~: $2q^2=p^2$. Donc $p^2$ est pair,
  ce qui implique que $p$ est pair (par contraposée du fait que le carré
  d'un impair est impair). Écrivons $p=2k$. Alors $2q^2=4k^2$, soit
  $q^2=2k^2$. Ainsi $q$ est pair. Mais alors $p$ et $q$ sont tous deux
  pairs, ce qui contredit $\pgcd(p,q)=1$.
\end{example}

\subsection{Raisonnement par récurrence}\label{subsec:recurrence}

\begin{remark}{Méthode}{rem:recurrence}
  Pour montrer qu'une propriété $P(n)$ est vraie pour tout $n\geq n_0$~:
  \begin{enumerate}[label=\textbf{(\arabic*)}]
    \item \textbf{Initialisation~:} Vérifier que $P(n_0)$ est vraie.
    \item \textbf{Hérédité~:} Montrer que pour tout $n\geq n_0$,
      $P(n)\implies P(n+1)$.
  \end{enumerate}
  On conclut par le \emph{principe de récurrence} que $P(n)$ est vraie
  pour tout $n\geq n_0$.
\end{remark}

\begin{example}{Somme des entiers}{ex:somme-entiers}
  \textbf{Énoncé.} Pour tout $n\geq 1$,\; $\displaystyle\sum_{k=1}^{n}k=\frac{n(n+1)}{2}$.

  \textbf{Preuve par récurrence.}

  \emph{Initialisation} ($n=1$)~: $\sum_{k=1}^{1}k=1=\frac{1\cdot 2}{2}$. \checkmark

  \emph{Hérédité.} Supposons la formule vraie au rang $n\geq 1$.
  Alors
  \[
    \sum_{k=1}^{n+1}k = \underbrace{\sum_{k=1}^{n}k}_{=\,n(n+1)/2} + (n+1)
    = \frac{n(n+1)}{2}+\frac{2(n+1)}{2}
    = \frac{(n+1)(n+2)}{2},
  \]
  ce qui est la formule au rang $n+1$.
\end{example}

\begin{remark}{Récurrence forte}{rem:recurrence-forte}
  Dans la \emph{récurrence forte} (ou \emph{complète}), l'hypothèse de
  récurrence porte sur \emph{tous} les rangs précédents~: on suppose
  $P(n_0),P(n_0+1),\ldots,P(n)$ pour en déduire $P(n+1)$. Cette
  variante est parfois indispensable (par exemple pour la factorisation
  en nombres premiers).
\end{remark}

% ============================================================
\section{Exercices}\label{sec:exo-ch0}
% ============================================================

\begin{exercise}{\basic{} Opérations ensemblistes}{exo:ens-ops}
  Soient $A=\set{1,2,3,4,5}$, $B=\set{3,4,5,6,7}$ et
  $E=\set{1,2,\ldots,10}$. Déterminer~:
  \begin{enumerate}[label=(\alph*)]
    \item $A\cup B$, $A\cap B$, $A\setminus B$, $B\setminus A$~;
    \item $A^c$, $B^c$, $(A\cup B)^c$, $(A\cap B)^c$~;
    \item Vérifier les lois de De Morgan (
ef{prop:prop:de-morgan}) sur cet
      exemple.
  \end{enumerate}
\end{exercise}

\begin{exercise}{\basic{} Injectivité et surjectivité}{exo:inj-surj}
  Pour chacune des applications suivantes, dire si elle est injective,
  surjective, bijective. Justifier.
  \begin{enumerate}[label=(\alph*)]
    \item $f\colon\Z\to\Z$, $n\mapsto 2n$.
    \item $g\colon\Z\to\Z$, $n\mapsto n+3$.
    \item $h\colon\R\to\R$, $x\mapsto x^3-x$.
    \item $\vphi\colon\R\to\R_+$, $x\mapsto e^x$.
  \end{enumerate}
\end{exercise}

\begin{exercise}{\basic{} Tables de $\F_7$}{exo:F7}
  \begin{enumerate}[label=(\alph*)]
    \item Construire la table de multiplication de $\F_7^*=\set{1,2,3,4,5,6}$.
    \item En déduire l'inverse de chaque élément non nul.
    \item Résoudre l'équation $3x+5=2$ dans $\F_7$.
  \end{enumerate}
\end{exercise}

\begin{exercise}{\intermediate{} Composition de bijections}{exo:comp-bij}
  Soient $f\colon E\to F$ et $g\colon F\to G$ deux bijections. Montrer
  que $g\circ f$ est une bijection et que $(g\circ f)^{-1}=f^{-1}\circ g^{-1}$.
\end{exercise}

\begin{exercise}{\intermediate{} Sous-groupes de $(\Z,+)$}{exo:sous-groupes-Z}
  Montrer que les sous-groupes de $(\Z,+)$ sont exactement les ensembles
  de la forme $n\Z\deff\setcond{nk}{k\in\Z}$, pour $n\in\N$.

  \emph{Indication~:} utiliser la division euclidienne.
\end{exercise}

\begin{exercise}{\intermediate{} Récurrence --- inégalité de Bernoulli}{exo:bernoulli}
  Montrer par récurrence que pour tout $x\geq -1$ et tout $n\in\N$~:
  \[
    (1+x)^n \geq 1+nx.
  \]
\end{exercise}

\begin{exercise}{\intermediate{} Corps fini $\F_2$ et parité}{exo:F2}
  Dans le corps $\F_2=\set{0,1}$, on a $1+1=0$.
  \begin{enumerate}[label=(\alph*)]
    \item Écrire les tables d'addition et de multiplication de $\F_2$.
    \item Résoudre le système linéaire suivant sur $\F_2$~:
      \[
        \begin{cases}
          x+y+z=1,\\
          x+z=0,\\
          y+z=1.
        \end{cases}
      \]
    \item Interpréter la résolution en termes de parité.
  \end{enumerate}
\end{exercise}

\begin{exercise}{\challenging{} $\Q(\sqrt{2})$ est un corps}{exo:Q-sqrt2}
  Soit $\Q(\sqrt{2})\deff\setcond{a+b\sqrt{2}}{a,b\in\Q}$.
  \begin{enumerate}[label=(\alph*)]
    \item Montrer que $\Q(\sqrt{2})$ est stable par addition et
      multiplication.
    \item Montrer que tout élément non nul de $\Q(\sqrt{2})$ admet un
      inverse dans $\Q(\sqrt{2})$.

      \emph{Indication~:} montrer que si $a+b\sqrt{2}\neq 0$ avec
      $a,b\in\Q$, alors $a^2-2b^2\neq 0$, puis conjuguer.
    \item En déduire que $(\Q(\sqrt{2}),+,\times)$ est un corps. Quelle
      est sa caractéristique~?
  \end{enumerate}
\end{exercise}

\begin{exercise}{\challenging{} Automorphismes de groupes}{exo:aut-groupe}
  Soit $(G,\star)$ un groupe. Un \emph{automorphisme} de $G$ est une
  bijection $\vphi\colon G\to G$ telle que
  $\vphi(a\star b)=\vphi(a)\star\vphi(b)$ pour tous $a,b\in G$.
  \begin{enumerate}[label=(\alph*)]
    \item Montrer que l'ensemble $\Aut(G)$ des automorphismes de $G$,
      muni de la composition, est un groupe.
    \item Déterminer $\Aut(\Z/n\Z,+)$ pour $n=2,3,4$.
  \end{enumerate}
\end{exercise}

% ============================================================
\section*{Résumé du chapitre}
\addcontentsline{toc}{section}{Résumé du chapitre}
% ============================================================

\begin{framed}
\textbf{Ce qu'il faut retenir.}

\begin{enumerate}[label=\textbf{\arabic*.},leftmargin=2em]
  \item \textbf{Ensembles.} Maîtriser les opérations $\cup$, $\cap$,
    $\setminus$, $\times$, $\mathcal{P}(E)$, ainsi que les lois de De Morgan.

  \item \textbf{Applications.} Savoir caractériser injectivité,
    surjectivité et bijectivité. Savoir composer des applications et
    utiliser les propriétés de la composition.

  \item \textbf{Relations.} Distinguer relation d'équivalence (classes,
    quotient) et relation d'ordre (total ou partiel).

  \item \textbf{Structures algébriques.}
    \begin{itemize}
      \item \emph{Groupe}~: un ensemble avec une loi associative, un neutre
        et des symétriques.
      \item \emph{Anneau}~: deux lois (addition formant un groupe abélien,
        multiplication associative avec unité, distributivité).
      \item \emph{Corps}~: un anneau commutatif où tout élément non nul
        est inversible.
    \end{itemize}

  \item \textbf{Exemples de corps.} $\Q$, $\R$, $\C$ (caractéristique~$0$),
    $\F_p$ pour $p$ premier (caractéristique~$p$). En algèbre linéaire,
    $\K$ désigne le corps de base.

  \item \textbf{Techniques de démonstration.}
    \begin{itemize}
      \item \emph{Preuve directe}~: supposer $P$, déduire $Q$.
      \item \emph{Contraposée}~: montrer $\lnot Q\implies\lnot P$.
      \item \emph{Absurde}~: supposer $\lnot P$, obtenir une contradiction.
      \item \emph{Récurrence}~: initialisation + hérédité.
    \end{itemize}
\end{enumerate}
\end{framed}
